\section{Modernized Nuclotron}
\label{sec:nuclotron}

\subsection{Present Nuclotron lattice}

Figure~\ref{fig:nica_complex} shows the NICA accelerator complex, which
comprises the Booster, the Nuclotron and the Collider. The Booster and the
Nuclotron serve as the injector of heavy ions, polarized protons and
deuterons into the collider. The Nuclotron lattice was originally designed to
be as close to circular as possible and has eight superperiods.
Figure~\ref{fig:nuclotron_twiss} shows the Twiss functions of one superperiod,
which consists of a bending arc of three FODO cells, each with four bending
magnets with a maximum field of \SI{1.8}{T}, and a straight section of one FODO
cell. The Nuclotron circumference is \SI{251}{m}.

\begin{figure}[htbp]
  \centering
  \includegraphics[width=0.7\textwidth]{nica_complex}
  \caption{NICA accelerator complex.}
  \label{fig:nica_complex}
\end{figure}

\begin{figure}[htbp]
  \centering
  \includegraphics[width=0.7\textwidth]{nuclotron_twiss}
  \caption{Twiss functions of one Nuclotron superperiod.}
  \label{fig:nuclotron_twiss}
\end{figure}

The main drawback of the Nuclotron, both as a collider injector and as a
stand-alone ring, is the limited length of the drift spaces in which the
equipment for these two functions must be installed. All modernization options below
are assessed with respect to this limitation.

\subsection{Modernization options}

Two modernization options of the Nuclotron have been studied.

The first option, Nuclotron~2.1, is a lattice that preserves all properties of
the existing accelerator and allows deuteron EDM experiments with electric
deflectors or Wien filters installed in the ring. Either device keeps the spin
of the polarized beam along its direction of motion around the ring. Within the
eight-superperiod lattice, $N_{super} = 8$, this requires implementing the
frozen (or quasi-frozen) spin concept by lengthening the straight sections
between the arcs, ensuring zero dispersion in the straight sections, and, if
possible, keeping the ring circumference so that the equipment fits into the
existing tunnel. An axion search could also be included in this framework,
although the treatment of the systematic errors of such an experiment is not
yet fully understood. In the Nuclotron~2.1 option, all key collider
parameters, such as the transition energy and the electron cooler parameters,
remain unchanged. This creates a serious problem at the interface between the
collider and the Nuclotron: crossing the transition energy with a cooled beam
leads to irreversible beam degradation and a lower final luminosity.

The Nuclotron~2.2 option introduces modifications that ensure compatibility
with the collider. The injection energy of polarized protons into the collider
remains at 2--\SI{3}{GeV}. To eliminate the effect of transition crossing after
electron cooling, the transition energy of the collider is to be raised above
its maximum operating energy. In addition, a lower maximum energy of heavy ions
and polarized particles in Nuclotron~2.2 is considered. This frees space in the
ring for Wien filters and extends the capabilities of the Nuclotron for EDM and
axion experiments. A lower Nuclotron energy does not degrade the complex as a
whole, since the collider can take over some functions of the Nuclotron and
provide proton and heavy-ion beams over the entire Nuclotron energy range.

\subsection{Nuclotron 2.1: quasi-frozen spin for deuterons}

Adapting the Nuclotron lattice to the measurement of the deuteron EDM involves
three problems: lengthening the straight sections, suppressing the dispersion
in the straight sections, and keeping the spin direction along the momentum
around the ring.

The first problem is solved by raising the maximum field of the bending
magnets to \SI{1.8}{T} and by merging the two adjacent magnets between the
quadrupoles into one magnet. At the same time, the dispersion is suppressed by
choosing the horizontal phase advance in the arcs. One superperiod of the
modernized Nuclotron~2.1 is shown in Fig.~\ref{fig:nuclotron21}. The main
parameters of Nuclotron~2.1 are: circumference \SI{251}{m}; betatron tunes
$Q_x = 9.78548$, $Q_y = 10.6839$; momentum compaction factor 0.0134394; slip
factor $-0.0222388$; natural chromaticity $-16.0394$ (horizontal) and
$-17.8984$ (vertical). In the figure, D denotes a defocusing quadrupole, F a
focusing quadrupole, B a dipole magnet and ED an electrostatic deflector. The
parameters of all options are summarized in Table~\ref{tab:nuclotron}.

\begin{figure}[htbp]
  \centering
  \includegraphics[width=0.8\textwidth]{nuclotron21_twiss}
  \caption{Twiss functions of one superperiod of the modernized Nuclotron~2.1
    with two electrostatic deflectors.}
  \label{fig:nuclotron21}
\end{figure}

With the bending field raised to \SI{1.8}{T}, the length of the straight
section in each superperiod increases from 7.3 to \SI{10.5}{m}. The third
problem, keeping the spin direction relative to the momentum as required for
the detection of the deuteron EDM signal, is solved with negative-curvature
electrostatic deflectors in each superperiod (Fig.~\ref{fig:nuclotron_ed}).

In the quasi-frozen spin concept, the magnetic and electric fields act in
different elements. The spin direction therefore oscillates relative to the
direction of motion, since the spin rotates in opposite directions in the
magnetic and electrostatic bending elements. The deuteron magnetic anomaly is
small, $G = -0.143$. The spin oscillates around the momentum in each magnetic
arc within half of the spin phase advance,
$\Phi_s = \frac{1}{2N}\,2\pi\gamma G$, and each time it is turned back by
$-\Phi_s$ in the electrostatic deflector; here $N$ is the arc superperiodicity.
This keeps the average spin direction along the momentum around the whole ring,
in accordance with the quasi-frozen spin concept~\cite{Senichev:2020qfs}.

For the Nuclotron, $N = 8$. Owing to the small value of $\Phi_s$, the effective
EDM signal is reduced only by the factor
\begin{equation}
  J_0(\Phi_s) \approx 1 - \frac{\Phi_s^2}{4},
\end{equation}
i.e., by a few percent at most.

To determine the parameters of the electrostatic deflectors in the
quasi-frozen spin mode for deuterons, consider the lattice as consisting of
two different parts (Fig.~\ref{fig:nuclotron_ed}). The magnetic arcs bend the
beam by $\Phi^B = 2\pi/N + 2\alpha$ each and rotate the spin in the horizontal
plane relative to the momentum by $\Phi^B_s = \nu^B_s\Phi^B$. The electrostatic
arcs with negative-curvature deflectors bend the beam by $\Phi^E = -2\alpha$
each and rotate the spin relative to the momentum in the opposite direction by
$\Phi^E_s = \nu^E_s\Phi^E$. The quasi-frozen spin condition requires
$\Phi^E_s = -\Phi^B_s$. Since in an electrostatic deflector the spin rotates
relative to the momentum many times faster than in a magnetic field, the two
types of arcs are related by
\begin{equation}
  \nu_s^B\left(\frac{2\pi}{N}+2\alpha\right) = \nu_s^E\, 2\alpha,
  \qquad
  \alpha = \frac{1}{\nu_s^E/\nu_s^B-1}\,\frac{\pi}{N} = -\gamma^2\frac{G}{G+1}\,\frac{\pi}{N},
  \label{eq:nuclotron_alpha}
\end{equation}
where $\nu_s^E$ and $\nu_s^B$ are given by Eqs.~\eqref{eq:nu_e}
and~\eqref{eq:nu_b}. For given bending magnets and a deflector field
$E_{ED} = \SI{130}{kV/cm}$, one obtains $\nu_s^E/\nu_s^B$, the parameter
$\alpha$ and the radius of curvature of the deflector $R_{ED} = \kappa/E_{ED}$,
where $\kappa = p_0\beta c/e$ is the electric rigidity and
$p_0 = \gamma m\beta c$, so that
$R_{ED} = \frac{mc^2}{e}\frac{\gamma\beta^2}{E_{ED}}$. The deflector length is
then
$L_{ED} = 2\alpha R_{ED} = -2\frac{\gamma^2 G}{G+1}\frac{\pi}{N}R_{ED}$.
Figure~\ref{fig:nuclotron_ed} shows how the electrostatic deflectors are
inserted to realize the quasi-frozen spin lattice of the Nuclotron. At the
optimal polarimeter energy $W = \SI{270}{MeV}$~\cite{Kurilkin:2010ym},
$\alpha = 0.026\pi$ and the required length of the electrostatic channel is
\SI{7.3}{m}. The EDM signal relative to the frozen spin case,
$S^{QFS}_{EDM}/S^{FS}_{EDM} = 1-\alpha^2/4 \approx 0.998$, is reduced by only
0.2\%.

The electrostatic arc with negative curvature leaves room for the necessary
equipment in the straight section between two magnetic arcs; when the
electrostatic deflectors are switched off, the beam passes straight through
this section.

\begin{figure}[htbp]
  \centering
  \includegraphics[width=0.8\textwidth]{nuclotron_ed_insert}
  \caption{Electrostatic insertion compensating the spin rotation in the
    magnetic arc.}
  \label{fig:nuclotron_ed}
\end{figure}

A lattice with twice the number of superperiods, 16, was also considered.
By splitting the double bending magnets into individual magnets, the number of
FODO cells per superperiod is doubled, and with a \ang{60} phase advance per
FODO cell the dispersion is suppressed in every second straight section. In
this version $\alpha = 0.013\pi$, and the EDM signal relative to the frozen spin
case is $1-\alpha^2/4 \approx 0.9996$, i.e., reduced by only 0.04\%, which makes
this lattice practically equivalent to the frozen spin one. This lattice can
also be used for the proton EDM, since the spin rotation angle of protons in
the magnetic arcs is only about \ang{30}.

\subsection{Nuclotron 2.2: quasi-frozen spin with Wien filters for deuterons}

In the Nuclotron~2.2 lattice, a small vertical magnetic field of about
\SI{100}{mT} is added to compensate the Lorentz force of the electric field,
so that the electrostatic deflectors become straight. With the magnetic field
added, the E+B element is a Wien filter. The ring then consists of arcs
connected by straight sections with Wien filters, in which crossed electric
and magnetic fields keep the particles on a straight path (zero Lorentz force)
and compensate the spin rotation accumulated in the magnetic arcs. The total
length of the Wien filters is determined by the beam parameters and by the
required compensation of the spin rotation in the arcs, whereas their number
is chosen for convenient placement around the ring.

With $N$ magnetic arcs in total, each arc bends the particles by
$\Phi^B_{arc} = 2\pi/N$ and rotates the spin in the horizontal plane relative
to the momentum by $\Phi^{arc}_s = \gamma G\,\Phi^B_{arc}$. The Wien filters
rotate the spin relative to the momentum in the opposite direction. In the
electric field $E_{wf}$ the spin rotation angle is
\begin{equation}
  \Phi^E_{swf} = -\left(\gamma G+\frac{\gamma}{\gamma+1}\right)\beta^2\,\Phi^E_{wf},
\end{equation}
where $\Phi^E_{wf} = \frac{eE_{wf}}{m_0c^2\gamma\beta^2}L_{wf}$ is the momentum
rotation angle in the electric field of the Wien filter, and in the magnetic
field $B_{wf}$ it is
\begin{equation}
  \Phi^B_{swf} = (\gamma G+1)\,\Phi^B_{wf},
\end{equation}
where $\Phi^B_{wf} = \frac{eB_{wf}}{m_0c\gamma\beta}L_{wf}$ is the momentum
rotation angle in the magnetic field of a Wien filter of length $L_{wf}$.

Since the Lorentz force in the Wien filter vanishes, the angles must be equal,
$\Phi^E_{wf} = \Phi^B_{wf}$, and they can be expressed through the electric
field $E_{wf}$ and the length $L_{wf}$ of the filters installed in all $N$
straight sections. In the frozen spin option, the total spin precession
frequency relative to the momentum in the combined E+B bending element is
zero. The quasi-frozen spin option instead requires the total spin rotation
angle per turn in the E and B fields to vanish. The QFS condition therefore
reads
\begin{equation}
  (\gamma G+1)\,\Phi^B_{wf} - \left(\gamma G+\frac{\gamma}{\gamma+1}\right)\beta^2\,\Phi^E_{wf} = \gamma G\,\frac{2\pi}{N},
\end{equation}
which gives the parameters of one Wien filter:
\begin{equation}
  L_{wf}E_{wf} = \frac{2\pi}{N}\,\frac{G}{G+1}\,\frac{m_0c^2}{e}\,\beta^2\gamma^3,
  \qquad
  B_{wf} = -\frac{E_{wf}}{c\beta},
  \label{eq:nuclotron_wf}
\end{equation}
where $L_{wf}$ is the length of one Wien filter. The total length of all Wien
filters,
$L^\Sigma_{wf}E_{wf} = 2\pi\frac{G}{G+1}\frac{m_0c^2}{e}\beta^2\gamma^3$,
depends only on the particle species and energy.

Equations~\eqref{eq:nuclotron_alpha} and~\eqref{eq:nuclotron_wf}, derived for
electrostatic deflectors and for Wien filters, show that the type of the
compensating element does not matter. With the magnetic arc unchanged, the
Wien filters are shorter by the total length of the kickers, since a Wien filter
combines the functions of an electrostatic deflector and a kicker in one
element.

Figure~\ref{fig:nuclotron221} shows the Twiss functions of one superperiod of
the first version of the modernized lattice, Nuclotron~2.2.1. Two Wien filters
installed in the central straight sections of the superperiod restore the spin
direction after each magnetic half-arc. The spin deviation angle from the
momentum, $\alpha = \gamma G\,2\pi/N$, characterizes how much the quasi-frozen
spin lattice differs from the frozen spin one. In this lattice, the number of
arcs that determine the spin deviation equals the number of superperiods,
$N = 8$.

\begin{figure}[htbp]
  \centering
  \includegraphics[width=0.8\textwidth]{nuclotron221_twiss}
  \caption{Twiss functions of one superperiod of the modernized
    Nuclotron~2.2.1 with two Wien filters.}
  \label{fig:nuclotron221}
\end{figure}

For a maximum electric field of \SI{130}{kV/cm}, the magnetic field is below
\SI{80}{mT}. Such a low magnetic field simplifies the Wien filter design: a
permanent magnet or an air-core coil can be used. For these fields and a
deuteron energy of \SI{270}{MeV}, the total length of the Wien filters in all
straight sections is about \SI{52}{m}, and the length of one Wien filter is
\SI{3.25}{m}. At \SI{270}{MeV} (\SI{135}{MeV} per nucleon), the maximum spin
deviation is $\alpha\sim\pm\ang{3}$ over one arc, which is a good
approximation to a regular frozen spin lattice.

In the next version, Nuclotron~2.2.2, the Wien filters are distributed over
four straight sections of the superperiod (Fig.~\ref{fig:nuclotron222}). Each
arc of the superperiod is thus split into two arcs, and with respect to the
bending magnets the lattice has the superperiodicity $N = 16$, while the overall
superperiodicity of the lattice with all elements remains 8. This solves two
problems. First, the ring becomes a 16-sided polygon, i.e., closer to a circle,
and fits the existing tunnel better. Second, with a bending magnet length of
\SI{1.77}{m} and a Wien filter length of \SI{1.65}{m}, the spin deviation is
$\alpha\sim\pm\ang{1.5}$, half that of version~2.2.1. This brings the
properties of the quasi-frozen spin lattice much closer to those of the frozen
spin lattice.

\begin{figure}[htbp]
  \centering
  \includegraphics[width=0.8\textwidth]{nuclotron222_twiss}
  \caption{Twiss functions of one superperiod of the modernized
    Nuclotron~2.2.2 with four Wien filters.}
  \label{fig:nuclotron222}
\end{figure}

To increase the energy, the modification Nuclotron~2.2.3 with fewer empty
drift spaces and longer magnets was considered (Fig.~\ref{fig:nuclotron223}).
In this lattice, the maximum energy is \SI{3.5}{GeV} per nucleon for deuterons
and \SI{7.8}{GeV} for protons, which is much closer to the present Nuclotron.

\begin{figure}[htbp]
  \centering
  \includegraphics[width=0.75\textwidth]{nuclotron223_twiss}
  \caption{Twiss functions of one superperiod of the modernized
    Nuclotron~2.2.3 with four Wien filters.}
  \label{fig:nuclotron223}
\end{figure}

\begin{table}[htbp]
  \centering
  \caption{Main parameters of the Nuclotron modifications.}
  \label{tab:nuclotron}
  \footnotesize
  \setlength{\tabcolsep}{2.5pt}
  \begin{tabular}{lcccccccc}
    \toprule
    \makecell[l]{Nuclotron\\lattice} & \makecell{Magnet\\length, m} & \makecell{Number of\\magnets}
      & \makecell{Total magnet\\length, m} & \makecell{$B_{max}$,\\T}
      & \makecell{$B_{max}R$,\\T$\cdot$m} & \makecell{$W_d^{max}$,\\GeV}
      & \makecell{$W_d^{max}$,\\GeV/u} & \makecell{$W_p^{max}$,\\GeV} \\
    \midrule
    Present & 1.44 & 96 & 138.24 & 1.8 & 39.6  & 10.14 & 5.07 & 10.97 \\
    2.1     & 2.35 & 48 & 112.77 & 1.8 & 32.2  & 8     & 4    & 8.79 \\
    2.2.1   & 1.78 & 48 & 85.3   & 1.8 & 24.4  & 5.7   & 2.85 & 6.47 \\
    2.2.2   & 1.78 & 48 & 85.44  & 1.8 & 24.5  & 5.7   & 2.85 & 6.47 \\
    2.2.3   & 2.10 & 48 & 101    & 1.8 & 28.93 & 7.0   & 3.5  & 7.8 \\
    \bottomrule
  \end{tabular}
\end{table}

\subsection{Nuclotron 2.2.2: quasi-frozen spin with Wien filters for protons}

With a magnetic superperiodicity of 16, the lattice can also be used for
protons at a reduced energy. The energy is then limited mainly by the free
space for the Wien filters, which is \SI{52}{m} in this lattice. The maximum
proton energy follows from Eq.~\eqref{eq:nuclotron_wf}:
\begin{equation}
  \gamma(\gamma^2-1) = L_{wf}\,\frac{eE_{wf}}{mc^2}\,\frac{N}{2\pi}\,\frac{G_p+1}{G_p},
\end{equation}
where $L_{wf}N = \SI{52}{m}$ and $G_p = 1.793$ is the proton magnetic anomaly.
Solving for $\gamma$ gives $\gamma = 1.08047$, i.e., $W_p\approx\SI{80}{MeV}$.
At this energy, the spin deviation in a magnetic arc is determined by
$\gamma G_p = 1.93711$, and the spin deviation angle is
$\alpha = \pm\frac{1}{2}\frac{\gamma G_p}{16}$, which amounts to
$\alpha = \ang{3.5}$.\todo[inline]{Check: with the factor $2\pi$ of
Eq.~\eqref{eq:nuclotron_wf}, $\alpha=\pm\frac12\gamma G_p\frac{2\pi}{16}\approx\ang{22}$.}
The same lattice as for deuterons can be used for protons, but at a lower
energy and with the Wien filters rotated by \ang{180}, since the magnetic
anomalies of the two particles have opposite signs. At a proton energy of
\SI{80}{MeV} the lattice thus behaves practically as a frozen spin lattice.
Figure~\ref{fig:pc_ay} shows the p--C analyzing power in terms of the
``average analyzing power'' taken from Ref.~\cite{Mcnaughton:1985wyt}.

\begin{figure}[htbp]
  \centering
  \includegraphics[width=0.6\textwidth]{pc_analyzing_power}
  \caption{Average p--C analyzing power $\overline{A_c}$ for
    $\theta_{lab} = \ang{5}$--\ang{20} at discrete energies (points) and its
    continuous energy dependence (solid line). The data are from SIN
    ($\triangle$), TRIUMF ($\times$) and LAMPF ($\bullet$)~\cite{Mcnaughton:1985wyt}.}
  \label{fig:pc_ay}
\end{figure}

Figure~\ref{fig:pc_ay} shows that the analyzing power near \SI{100}{MeV} is
about 0.2, i.e., 0.35 of its maximum value $A_c = 0.53$ at \SI{230}{MeV}. This
makes proton EDM studies at this energy in the Nuclotron~2.2.2 lattice
feasible in parallel with full-scale deuteron measurements.

Galactic axions can be searched for with the same equipment as the EDM.
Varying the spin precession frequency at a fixed beam energy turns the booster
with Wien filters into a broadband antenna that searches for axions through a
resonant full or partial spin flip in the oscillating pseudomagnetic field of
the axion halo of our Galaxy~\cite{Nikolaev:2022wyi}. The growth of the
horizontal polarization induced by axions is extracted by a Fourier analysis of
the oscillating vertical asymmetry in the internal
polarimeter~\cite{JEDI:2015vwa}.

\subsection{Maximum energy of heavy ions and protons in Nuclotron 2.2.2}

The calculations show that deuteron and proton EDM experiments in the
Nuclotron~2.2.2 lattice require free space for Wien filters with a total length
of \SI{52}{m}. With the tunnel circumference fixed at \SI{251}{m}, this is
possible either by raising the field of the bending magnets, which is already
set to \SI{1.8}{T}, or by reducing their total length. Assuming that the
bending field does not exceed \SI{1.8}{T}, the second option was chosen, which
inevitably lowers the maximum particle energy in the Nuclotron.

\subsection{Features of the modernized Nuclotron lattice}

The lattice of the modernized Nuclotron was designed with several goals.
First, the superperiod has a drift without bending magnets in its central
cell, where the dispersion function reaches its maximum and the orbit
curvature is zero ($\rho\to\infty$). This makes it possible to raise the
transition energy, since the momentum compaction factor is
\begin{equation}
  \alpha = \frac{1}{C}\oint\frac{D(s)}{\rho(s)}\,ds .
\end{equation}
Second, the number of superperiods $s = 8$ is only slightly larger than the
horizontal betatron tune $\nu_x$, which makes the momentum compaction factor
easy to adjust, as in the resonant lattices with a low or negative momentum
compaction factor~\cite{Senichev:1997kn,Ishi:1998tj,Papaphilippou:2010zz,Senichev:2008zza,Toelle:2007zz}:
\begin{equation}
  \alpha = \frac{\nu_x^3}{R}\sum_{k=0}^{\infty}\frac{a_k^2}{\nu_x^2-(ks)^2},
\end{equation}
where $a_k$ are the harmonics of the dispersion-generating perturbation.
Third, the central drift simplifies the placement of the three sextupole
families needed to control the betatron and spin chromaticities, whereas the
side drifts are convenient for the Wien filters. Finally, the maximum of the
dispersion at the center of the superperiod coincides with the minimum of the
horizontal beta function $\beta_x$.

\subsection{Summary}

If the Nuclotron is used as the booster of the collider, its final energy can
be lowered to 2--\SI{3}{GeV} without a significant degradation of the NICA
complex, because all physics experiments with polarized beams are transferred
to the collider. This choice has a further advantage: it shortens the total
length of the bending magnets, e.g., by reducing each double bending magnet of
the Nuclotron to one section, and frees about \SI{70}{m} of the orbit for
additional equipment, in particular for polarization-preserving solenoids. It
also provides room for EDM and axion search equipment directly in the
Nuclotron.

The material of this section has been reported
in~\cite{Senichev:2023vtx,Senichev:2024fey,Senichev:2025ude,Senichev:2026xca,Kolokolchikov:2025sew,Senichev:2025jqf}.
